% Source: https://artofproblemsolving.com/community/c3h1600256_2002_mathcounts_state_team_round_problem_5
% Stable identity: aops-topic-1600256-q01
All four wheels have the same size, so a wheel's rotation rate is proportional to its speed, and in one trip around the circle the distance it travels is proportional to the radius of its circle (both front wheels go around the same centre $O$ in the same time).

So the outside front wheel's circle has twice the radius of the inside one's. If the inside radius is $r$, the outside radius is $2r$, and the wheels are $2r-r=r=6$ feet apart:

\begin{asy}
usepackage("amsmath");
size(5.5cm);
defaultpen(fontsize(10pt));
draw(circle((0,0),6),blue+linewidth(1));draw(circle((0,0),12),red+linewidth(1));
dot((0,0));label("$O$",(0,0),SE);
draw((6,0)--(12,0),linewidth(1.2));
fill(box((5.6,-0.9),(6.4,0.9)),black);fill(box((11.6,-0.9),(12.4,0.9)),black);
label("$6$",(9,0),N);
draw((0,0)--6*dir(120),blue+dashed);label("$r$",3*dir(120),W,blue);
draw((0,0)--12*dir(215),red+dashed);label("$2r$",8*dir(215),N,red);
label("\mbox{}",(min(currentpicture,user=true)+max(currentpicture,user=true))/2,black); // keeps the text after the figure black when LaTeX typesets the labels inline
\end{asy}

The inside front wheel travels $2\pi\cdot6=\boxed{12\pi}$ feet.
